\documentclass[11pt]{article}
\usepackage[a4paper,margin=26mm]{geometry}
\usepackage{amsmath,amssymb,mathtools}
\usepackage{unicode-math}
\setmainfont{Latin Modern Roman}
\setmathfont{Latin Modern Math}
\usepackage{xurl,hyperref}
\hypersetup{hidelinks,pdftitle={UnsolvedMath Problem 10400012}}
\newcommand{\sref}[2]{\hyperlink{source:#1:#2}{[S#2]}}
\newcommand{\cataloguescope}[1]{#1}
\begin{document}
% UnsolvedMath ID 10400012; source catalogue code AMR-103-0012
\setcounter{section}{10400011}
\section{Problem 1.12 — Categorify other knot polynomials.}\label{10400012}
{\small\sffamily UnsolvedMath ID 10400012 · Topology}
\subsection{Definitions and mathematical statement}

\paragraph{Shared notation.}$\mathbb N=\{1,2,\ldots\}$, and $\mathbb Z,\mathbb Q,\mathbb R,\mathbb C$ denote the integers, rationals, reals and complex numbers. Any other conventions are specified locally.


A knot or link is a smooth oriented embedding of one or more circles in the oriented $3$-sphere, considered up to orientation-preserving ambient isotopy. Write $O$ for the unknot, which bounds an embedded disk. Categorifying a polynomial invariant means assigning graded homology groups, invariant under the relevant knot or link equivalence, whose graded Euler characteristic recovers that polynomial. For a multigrading $\mathbf j$ and homological degree $i$, this characteristic is $\sum_{i,\mathbf j}(-1)^i\dim H^{i,\mathbf j}\,\mathbf x^{\mathbf j}$, whenever the finiteness conditions make it defined.
\paragraph{Statement recorded in the supplied source.}
Construct meaningful categorifications of knot polynomials other than the Jones polynomial, including appropriate graded chain complexes, equivalence-invariant homology and a proof that the graded Euler characteristic is the target polynomial. Seek actual homological constructions with structural information, rather than assigning vector spaces solely from known coefficients. \sref{10400012}{2}
\subsection{Short English statement}
Construct invariant homology theories whose graded Euler characteristics recover other knot polynomials.
\subsection{Sources}
\begin{description}
\item[\hypertarget{source:10400012:1}{S1}] ULAM. AI and the UnsolvedMath Contributors, \emph{UnsolvedMath: A Curated Collection of Open Mathematics Problems}, record 10400012 (AMR-103-0012). CC BY 4.0. \url{https://huggingface.co/datasets/ulamai/UnsolvedMath}.
\item[\hypertarget{source:10400012:2}{S2}] T. Ohtsuki (editor), \emph{Problems on invariants of knots and 3-manifolds}, Geometry \& Topology Monographs 4 (2002), pp. 377--572 (published 2004), item 1.12 and its surrounding definitions and remarks. \url{https://arxiv.org/abs/math/0406190v1}.
\end{description}
\label{end:10400012}
\end{document}
